\author{Lothar}
\documentclass{mybib}


\begin{document}    
    {\hfill Sitzung vom 8. Nov 2026}

\section{Weihnachtsmarkt}

\subsection{Auszug aus dem Reglement}
\begin{itemize}
  \item Die Marktstände können am Donnerstag von 8.30 Uhr bis 16.30 Uhr eingerichtet werden. 
  \item \textbf{Offizielle Marktöffnungszeiten:}
  \begin{itemize}  
    \item Freitag von 17.00 Uhr bis 21.30 Uhr (22.00 Uhr für Food) 
    \item Samstag von 14.00 Uhr bis 21.30 Uhr (22.00 Uhr für Food) 
    \item Sonntag von 10.30 Uhr bis 17.00 Uhr
  \end{itemize}
  \item Am Abend wird vom Veranstalter jeweils um 22.30 Uhr das Licht auf dem Marktareal gelöscht. 
  \item Der Aussteller verpflichtet sich, seinen Marktstand währen den offiziellen Marktöffnungszeiten geöffnet zu halten. 
  \item Werden die Marktöffnungszeiten nicht eingehalten, wird der Aussteller nur einmal Verwarnt. Bei erneuter Widerhandlung wird der Stand geschlossen. Bei mehrmaligen Verfehlungen wird der Austeller für den Adventsmarkt auf immer ausgeschlossen.
\end{itemize}

\subsection{Aufstellen:}
\begin{itemize}

  \item Die Marktstände können am Eröffnungstag ab 8:30 Uhr bezogen werden und müssen bis 16:30 Uhr eingerichtet sein.
  \item Der Bezug des Marktstandes darf erst nach Anmeldung beim OK erfolgen.
  \item Der Aussteller ist verpflichtet den Marktstand weihnachtlich zu dekorieren, ohne die vom Veranstalter montierte Beleuchtung zu verdecken und das Dach zu beschädigen (Schrauben, Nägel, …).
  \item Der Veranstalter erstellt für jeden Marktstand die elektrische Grundversorgung. Bei Einzelnen besteht eine elektrische Versorgung.
  \item Abfälle sind vom Aussteller zu entsorgen. -- Während der offiziellen Marktöffnungszeiten dürfen keine vollen Abfallsäcke sichtbar herumstehen. -- Jeden Abend hat der Aussteller die herumliegenden Abfälle einzusammeln.
\end{itemize}

\subsection{Arbeiten:}
\begin{enumerate}
  \item Bereitstellung Material bis Mittwoch 25.11 im Lokal
  \item Bereitstellung der Dekoration bis Mittwoch 25.11
  \item Aufstellen und Einrichten des Standes ab Donnerstag 8:30Uhr möglich
  \item Besetzung des Standes während den Öffnungszeiten
  \item Verpflegung und Ablösung
  \item …
\end{enumerate}
\subsection{Einsatz:}
Ziel ist es, dass der Stand durchgehend min. 2 Personen besetzt ist. Optimal 3 Personen.
\subsubsection{Marie-Therese und Armin:}
Do und Fr von 17:00Uhr - 20:00Uhr. Samstag Frei und Sonntag nach Bedarf.
\subsubsection{Ueli:}
Samstag von 17:00Uhr bis 20:00Uhr Sonntag 14:00Uhr bis 17:00Uhr und aufräumen.
\subsubsection{Ursula:}
Donnerstag und Freitag von 20:00Uhr bis Schluss. Samstag: von 20:00Uhr bis Schluss
\subsubsection{Janina und Lothar:}
Einer von beiden ist immer da.
\subsubsection{Josette:}
Hilft beim Aufstelle, Abbauen reinigen und Verpflegung. Mädchen für alles.
\\
\\
Marie-Therese schaut sich zusammen mit Josette die Deko an. Aufstellen und Dekorieren Donnerstag ab 15:00Uhr. Material hinfahren ausladen und wenn das Auto weg ist hat man Zeit.\\
\textit{Wichtig! Wir dürfen keine Esswaren oder Getränke verkaufen. Dazu braucht es eine besondere Genehmigung. Selber Essen und Trinken ist OK.}

\subsection{Material:}
Janina und Lothar werden Bücher, Flyer und das auszustellende Material besorgen.

\textbf{Auswahl vom Bücher aus dem MNR Verlag:}
\begin{enumerate}  
  \item MNR Zeitschriften
  \item Israel Kalender
  \item Weihnachtsliteratur
  \item Gemalte Steine von Ursula
  \item Gratis Bibeln
  \item Gratis Kalender
  \item Gratis Flyer
\end{enumerate}

\section{Tag 1 Donnerstag}
\subsection{Einteilung}
\begin{center}
\begin{tabular}{ c | l l }
  Zeit & 17:00 & 21:00\\
  Donnerstag & J-U-L-M & J-U-L-M 
\end{tabular}
\end{center}
\subsection{Ablauf}
\begin{itemize}
  \item Besammlung um 14:00 Uhr in der Gemeinde
  \item Beschsprechung mit anschliessendem Gebet
  \item Einpacken der Literatur und Dekormaterial
  \item Mit dem Auto das Material zum Stand gefahren und aufgebaut
  \item Zurück zur Gemeinde
  \item Um 16:30 Uhr zum Markt.
  \item Da relativ wenig Platz können nur zwei den Tisch bedienen.
  \item Kalender wurden gut genommen, aber nur wenn man den Kalender den Leuten in die Finger drückt.
  \item Um 21:30 Uhr, Material mit einer Plane abgedeckt.
\end{itemize}
\section{Tag 2 Freitag}
\subsection{Einteilung}
\begin{center}
\begin{tabular}{ c | l l }
  Zeit & 17:00 - 19:00 & 19:00 - 21:30\\
  Freitag & L-M & J-U
\end{tabular}
\end{center}
\subsection{Ablauf}
\begin{itemize}
  \item Treffen Lothar und Marie-Therese 16:0n Uhr in der Gemeinde
  \item 17:00 Uhr öffnen des Standes
  \item 19:00 Uhr Schichtwechsel
\end{itemize}
\section{Tag 3 Samstag}
\subsection{Einteilung}
\begin{center}
\begin{tabular}{ c | l l l }
  Zeit & 14:00 - 17:00 & 17:00 - 19:00 & 19:00 - 21:30\\
  Samstag & Ue-Da & M-U & J-L
\end{tabular}
\end{center}
\subsection{Ablauf}
\begin{itemize}
  \item 
\end{itemize}
\section{Tag 4 Sonntag}
\subsection{Einteilung}

\begin{tabular}{ c | l l l }
  Zeit & 10:30 - 12:00 & 12:00 - 15:00 & 15:00 - 17:00\\
  Samstag & M-Da & L-J & Ue-U
\end{tabular}

\subsection{Ablauf}
\begin{itemize}
  \item 
\end{itemize}
\end{document}